\scanpage{90}

Position relative de deux racines (ou deux bases)

Soient $r,r'$ deux racines, correspondant à des bases \deleted{$\unsure{b,b'}$} $b,b'$

\[
\begin{aligned}
r'=r^{\unsure{*|+|\prime}}&\Longleftrightarrow r.r'=-2&\Longleftrightarrow b=b'\\
r'=-r&\Longleftrightarrow r.r'=2&\Longleftrightarrow b'\ \text{associé à }b\\
r'\ \text{orth. à }r&\Longleftrightarrow rr'=0&\Longleftrightarrow\operatorname{card}(b\cap b')=\underline{1}\\
\widehat{r,r'}=60^\circ&\Longleftrightarrow rr'=-1&\Longleftrightarrow\operatorname{card}(b\cap b')=3\ \deleted{\text{\unsure{et b,b′ associés}}}\\
\widehat{r.r'}=120^\circ&\Longleftrightarrow rr'=1&\Longleftrightarrow\operatorname{card}(b\cap b')=\deleted{\unsure{1|4}}0\\
&&\quad\deleted{\text{\unsure{et b,b′ associés}}}\ \inserted{\text{et }b,b'\text{ non \underline{associés}.}}
\end{aligned}
\]

Pour une racine $r$ donnée, il y a

\[
\begin{aligned}
1&\ \text{racine égale à }r,\ \text{soit }r\\
1&\ \text{racine opposée à }r,\ \text{soit }-r\\
30&\ \text{racines orthogonales à }r\\
20&\ \deleted{\text{\unsure{racines}}}\ \text{faisant un angle de }60^\circ\ \text{avec }r\\
20&\quad\rule{6em}{.3pt}\quad120^\circ\quad\rule{1.5em}{.3pt}\quad r
\end{aligned}
\]

\[
\begin{aligned}
\deleted{\text{\unsure{couples}}}\text{\unsure{paires|couples} de racines égales}&\quad72\\
\rule{5em}{.3pt}\ \text{opposées}&\quad72/2=36\\
\deleted{\text{\unsure{couples}}}\text{\unsure{paires|couples} de racines orthogonales}&\quad72\cdot15=1080\\
\deleted{\text{\unsure{couples}}}\text{\unsure{paires|couples} de racines à }60^\circ&\quad72\cdot10=720\\
\rule{4em}{.3pt}\quad120^\circ&\quad\underline{\underline{72\cdot10=720}}
\end{aligned}
\]

plans de \inserted{paires de} racines non orthogonales (= plans spéciaux) \quad $720:6=120\deleted{\unsure{=2^3\cdot3}}$

% archive sheet number: 89
